\subsection{I-21. A checked width-three specialization of the rank-preservation construction}
\label{I-21}

\paragraph{Source and status.}
The mechanism below is a specialization of Liu (2026),
\emph{Black-box identity testing for affine orbits of read-once
oblivious algebraic branching programs}, Hexagon
\texttt{2611.00634v1}, Section~3.1, Theorem~3.2 and Lemma~3.3,
and Section~3.2, Lemmas~3.4 and~3.6, printed pp.~5--7.
This is an unrefereed preprint dated November 19, 2026; its host
discloses primarily AI-generated text and no human-understood parts.
Its identifier is a Hexagon identifier, not an arXiv identifier.
We give a complete proof of the narrower original-coordinate,
width-three statement.  The affine-orbit assertions, the general-width
bounds, and the other sections of that manuscript are not being
certified here.  No historical priority is claimed.

\paragraph{Statement.}
Let $F$ be any field, let $n,D\ge1$, and let $f\in F[x_1,\ldots,x_n]$
have total degree at most $D$.  Suppose that for some permutation
$\sigma$ it has a constant-boundary representation
\[
 f=uA_1(x_{\sigma(1)})\cdots A_n(x_{\sigma(n)})v,
 \qquad A_i\in F[x_{\sigma(i)}]^{3\times3}.
\]
Smaller widths and rectangular layers are included by zero padding.
Let $E/F$ be a working field with at least
\[
 M=\max\{n,D+1,(n-1)D+1\}
\]
distinct available elements.  Choose distinct $a_1,\ldots,a_n\in E$,
and let $\lambda_i(Y)$ be their degree-$(n-1)$ Lagrange polynomials.
Define the twenty-one-parameter map
\[
 \mathcal G_i(t,\mathbf y,\mathbf z)
   =t+\sum_{j=1}^{10}z_j\lambda_i(y_j),
 \qquad 1\le i\le n.
 \tag{I.22}\label{eq:I-width3-generator}
\]
Then $f\ne0$ implies $f(\mathcal G)\ne0$ as a polynomial over $E$.
No variable order is an input to the map.

Put $R=(n-1)D+1$.  Take any sets $T,Y\subset E$ with
$|T|=D+1$ and $|Y|=R$.  For distinct
$\theta_0,\ldots,\theta_D\in E$, put
\[
 \Sigma_{10,D}=
 \{(\theta_{b_1},\ldots,\theta_{b_{10}}):
       b_j\ge0,\ \sum_j b_j\le D\}.
\]
The evaluations of $\mathcal G$ on
$T\times Y^{10}\times\Sigma_{10,D}$ form a hitting set of size at most
\[
 (D+1)\bigl((n-1)D+1\bigr)^{10}\binom{D+10}{10}.
 \tag{I.23}\label{eq:I-width3-size}
\]
This is polynomial in $n,D$ and is valid in every characteristic.
The field convention is explicit: if $E\ne F$, the black box must
allow evaluation over $E$.  No base-field query simulation is being
assumed.  In particular this does not contradict the finite-field
obstruction in I-16.

\paragraph{1. The coefficient-space lemma needed at width three.}
For a field $L$, work in
\[
 \mathcal R=L[Y_1,\ldots,Y_m]/(Y_1^{e_1+1},\ldots,Y_m^{e_m+1}),
 \qquad e_i\ge0,
\]
and suppose $B_i$ depends only on $Y_i$.  A matrix has \emph{unit rank}
$r$ if its constant matrix has rank $r$ and all its $(r+1)$-minors
vanish in $\mathcal R$.  Unit rank zero means that the matrix is zero.
For any constant row $b$, if every interval of the word
$B_1\cdots B_m$ has unit rank in $\{0,2,3\}$, then
\[
 \operatorname{Coeff}(bB_1\cdots B_m)
   =\operatorname{span}_L
       \{[Y^\alpha](bB_1\cdots B_m):
                         |\operatorname{supp}\alpha|\le4\}.
 \tag{I.24}\label{eq:I-profile23}
\]
If every interval has unit rank three, the corresponding bound is two.
Here $\operatorname{Coeff}$ means the span of coefficient rows.

We prove this assertion, including the stable-projection step.
An interval of unit rank zero makes the whole product zero.  Otherwise
the full constant product has rank two or three.  In rank three all
constant factors are invertible.  Cut-dependent constant changes of
basis make them identities.  The coefficient row spaces of the
successive prefixes are then nested.  Starting with one row, they can
grow at most twice.  At a growth layer one new nonconstant local
coefficient increases support by at most one; at a stable layer
constant propagation preserves a spanning set.  This proves the
bound two, with actual coefficient rows, as in I-17.

In rank two write the constant matrices as $C_i$.  Choose a
two-dimensional row space $U_0$ on which right multiplication by
$C_1\cdots C_m$ is injective, and set $U_i=U_0C_1\cdots C_i$.
Choose a one-dimensional complement $V_m$ of $U_m$.  At cut $i$ take
\[
 V_i=\{z:zC_{i+1}\cdots C_m\in V_m\}.
\]
The suffix has rank at least two and maps $U_i$ isomorphically onto
$U_m$.  Thus $V_i$ has dimension one, complements $U_i$, and
$V_{i-1}C_i\subseteq V_i$.  Transport bases of the $U_i$ along the
constants and choose bases of the $V_i$.  The resulting constant
matrices are
\[
 C_i=\operatorname{diag}(I_2,q_i),\qquad q_i\in L.
 \tag{I.25}
\]
These basis changes preserve supports and all interval unit ranks.

Let $\pi$ denote projection of a row onto its first two coordinates.
For a starting constant row space $S_0$, propagate
$S_i=\operatorname{Coeff}(S_{i-1}B_i)$.  Since
$\pi(sC_i)=\pi(s)$, the spaces $\pi S_i$ are nested.
Consider first a segment on which they are all the same space $T$.
Choose basis lifts of $T$ in $S_0$, with initial cost zero, and a basis
of the kernel of $\pi|_{S_0}$, with initial cost one.  Define the cost
of an obtained coefficient row as its initial cost plus the number
of nonconstant local coefficients used.  We claim that cost at most
one spans the final space on this stable segment.

If every $q_i$ on the segment is nonzero, the constants are invertible.
All the spaces lie in $\pi^{-1}(T)$, of dimension $\dim T+1$.
If the initial kernel is present, the initial space is already this
entire ambient space and constant propagation suffices.  Otherwise
there is room for at most one dimension increase.  Before that
increase the propagated constant lifts span; at the increase one
nonconstant coefficient suffices; afterwards constant propagation
again spans.  The claimed cost bound follows.

If some $q_i=0$, let $c$ be the last such position, put
$P=B_c\cdots B_m$, and let $B$ be the matrix of the initial lifts
propagated through the constants before $c$.  The suffix $P$ has
unit rank two and constant matrix $\operatorname{diag}(I_2,0)$.
Write its upper blocks as $[U\ V]$.  The constant part of $U$ is
$I_2$, so $U$ is invertible in $\mathcal R$.  The vanishing
three-by-three determinant, or its Schur complement, shows that
every row $(x,y)$ in the row module of $P$ satisfies
$y=xU^{-1}V$.

Let $H=\pi B$, whose rows form a basis of $T$.  If $T=0$, this row
graph forces every propagated row through $P$ to be zero.  Otherwise
write
\[
 \pi(BP)=QH,\qquad \pi(sP)=\eta H\quad(s\in S_{c-1}).
\]
Such expressions exist because all coefficient projections stay in
$T$, and $Q(0)=I$.  Therefore $Q$ is invertible.  The row graph gives
\[
 sP=\eta Q^{-1}BP.
 \tag{I.26}
\]
Coefficients of a polynomial matrix times $BP$ are linear
combinations of coefficient rows of $BP$, also in the truncated
ring.  Conversely the rows of $B$ belong to $S_{c-1}$.  It follows
that the final coefficient space is exactly
$\operatorname{Coeff}(BP)$.
Restart just after $c$ from $\operatorname{Coeff}(BB_c)$.
Constant lifts have cost zero and a kernel basis has cost at most
one, since only one variable position was used.  The remaining
constants are invertible; the preceding case therefore proves the
same cost bound.  This completes the stable-segment claim.

For the original single starting row $b$, first suppose $\pi b\ne0$.
The active projected space starts in dimension one and has at most
one strict growth.  The first stable segment costs at most one,
the growth factor at most one more, and the final stable segment
at most one more: support at most three suffices.  If $\pi b=0$,
the first stable segment has zero projected space.  The
stable-segment claim, with the sole initial row assigned cost one,
says that this segment is spanned by coefficients with no
nonconstant positions, or is zero.  There can subsequently be at
most two projection-growth factors.  Each costs at most one, and
each following stable segment at most one.  Support at most four
therefore suffices.  At every restart, lifts and kernel vectors
are linear combinations of the coefficient rows already retained;
the span filtration by support is preserved.  This proves (I.24).

The bound four in (I.24) is sharp.  With $P=\operatorname{diag}(1,1,0)$,
take
\[
 B_1=P+Y_1E_{31},\quad B_2=I+Y_2E_{13},\quad
 B_3=P+Y_3E_{32},\quad B_4=I+Y_4E_{23},\quad b=e_3^T.
\]
Every interval containing $B_1$ or $B_3$ has generic and constant
rank two, and the other intervals are invertible.  The output row is
\[
 Y_1e_1^T+Y_1Y_2Y_3e_2^T+Y_1Y_2Y_3Y_4e_3^T.
\]
Its third coefficient direction requires support four.

\paragraph{2. Lowest homogeneous components.}
Suppose every interval of a polynomial matrix word has the same
generic rank as at zero, with these ranks in $\{2,3\}$.  For constant
endpoints $b,c$, every monomial in the lowest nonzero homogeneous
component of $bB_1\cdots B_mc$ has support at most four.  If all
interval ranks are three, the bound is two.

Indeed, choose such a monomial $Y^\alpha$ and reduce modulo
$Y_i^{\alpha_i+1}$ for every $i$.  Only monomials dividing
$Y^\alpha$ survive.  Lowest-degree minimality makes the surviving
scalar polynomial exactly its nonzero $Y^\alpha$ term.  All the
intervals retain their unit ranks: the pivot constant minor remains
invertible and the higher minors were already zero polynomially.
Apply (I.24), or its invertible version.  If the support of $\alpha$
were too large, the final boundary $c$ would annihilate every
spanning coefficient row and hence also the surviving one, a
contradiction.

\paragraph{3. A finite collection of local rank tests.}
Augment the original matrix word by the constant factors $e_1u$
and $ve_1^T$.  Its full product has only one possibly nonzero entry,
namely $f$ in position $(1,1)$.  All intervals have a well-defined
generic rank over the rational function field.  Translation of
each distinct variable by any center in an overfield preserves
these generic ranks.

Fix $r=1$ or $r=2$, and suppose all intervals of generic rank
greater than $r$ are already preserved at the current center.
For any target interval $[l,h]$ of rank $r$, let $[l,i_0]$ be its
first prefix of rank $r$ and require this prefix to be preserved.
For each $i_0\le i<h$, choose the shortest rank-$r$ suffix $[j,i]$
of $[l,i]$, and require $[j,i+1]$ to be preserved.  Each such test
has generic rank $r$.  Its interior $[j+1,i]$, and the prefix
$[l,i_0-1]$ of the first test, are empty or have generic rank
greater than $r$ by minimality.  Thus a test has one such chunk
and at most one factor at each end.

These tests suffice.  Once $[l,i]$ is preserved, its specialized
row space is contained in that of $[j,i]$.  Both have dimension
$r$: the first has rank $r$, and the second cannot exceed its
generic rank $r$.  Their row spaces are equal.  Multiplying by
the next specialized factor and using the tested rank of
$[j,i+1]$ proves preservation of $[l,i+1]$.  This proves the
assertion by induction, for all target intervals.  Select one
nonzero $r$-minor of each test; no efficient selection or knowledge
of the hidden word is needed in the black-box construction.

For $r=2$, each selected minor is an entry of the word of second
compound matrices, by Cauchy--Binet.  These compounds have size
three.  The interior chunk and every one of its subintervals have
rank three, and so do their constants.  In a lowest-degree
monomial of the selected minor, extract the prescribed powers at
the at most two end factors.  The remaining scalar chunk
contraction is still in its lowest degree.  Part~2 gives support
at most two inside the chunk.  The whole monomial consequently
has support at most four.

For $r=1$ the compounds are the original matrices.  The chunk and
its subintervals have ranks in $\{2,3\}$, already preserved at
the center.  The same argument gives support at most four
inside the chunk and at most six in the whole selected entry.
Empty chunks and constant end factors only reduce these bounds.

\paragraph{4. Simultaneous preservation with four then six pairs.}
For $k\ge1$ write
\[
 \operatorname{SV}_{k,i}(\mathbf y,\mathbf z)
      =\sum_{j=1}^kz_j\lambda_i(y_j).
\]
Over any coefficient overfield of $E$, this map detects any
nonzero polynomial having a nonzero monomial of support at most
$k$: set the selectors to the Lagrange nodes of that support,
and the unused amplitudes to zero.  The resulting restriction
to that coordinate subspace retains the chosen coefficient.
In particular it detects a nonzero homogeneous polynomial all
of whose monomials have support at most $k$.

Start at the common center $(t,\ldots,t)$ over $E(t)$.
An interval of generic rank three consists entirely of invertible
univariate factors.  Its determinant at the center is a product
of nonzero univariate determinant polynomials evaluated at the
transcendental $t$, and is therefore nonzero.  Thus all rank-three
intervals are preserved simultaneously.

Add an $\operatorname{SV}_4$ perturbation with fresh independent
selectors and amplitudes.  For every chosen rank-two test minor,
its expansion about the current center is a nonzero polynomial
whose lowest homogeneous component $H$ has support at most four
in each monomial, by part~3.  Hence $H(\operatorname{SV}_4)\ne0$.
The map is homogeneous of degree one in its amplitudes.  The
lowest amplitude degree in the full substituted minor is
therefore exactly the degree of $H$, and its component is
$H(\operatorname{SV}_4)$; higher homogeneous terms cannot cancel
it.  Each selected test minor is consequently nonzero in the
same rational function field of the new parameters.  This is
simultaneous preservation, not a separate numeric choice for
each test: finitely many nonzero polynomials have nonzero product
in the parameter polynomial ring.  Earlier preserved minors
remain nonzero because setting the new amplitudes to zero
recovers their previous values.  The local-test induction now
preserves every interval of ranks two and three.

Over that enlarged rational function field, add an independent
$\operatorname{SV}_6$ perturbation.  The rank-one test components
have support at most six by part~3, so the same argument preserves
all these tests and all earlier ranks.  The augmented full word
has generic rank one whenever $f\ne0$.  Its specialized rank is
therefore one, which says precisely that the resulting scalar is
nonzero.  The two perturbations use disjoint parameter pairs;
their sum is $\operatorname{SV}_{10}$.  This proves (I.22).
All uses of the variable order occurred only in this existential
proof about the hidden representation.  The actual map is the
same for every order.

For comparison with the failed repair I-18, the polynomial
$\prod_{j=1}^m(x_{2j}-x_{2j-1})$ already survives one interpolation
pair: its pullback is
$z^m\prod_j(\lambda_{2j}(y)-\lambda_{2j-1}(y))$, which is nonzero
because the distinct Lagrange basis polynomials have nonzero
differences.  Interpolation directions are generally dense;
the construction therefore does not require a sparse perturbation
point to witness every nonzero polynomial.

\paragraph{5. Finite sampling and complexity.}
Since $f$ has total degree at most $D$, its pullback under (I.22)
has degree at most $D$ in $t$, at most $(n-1)D$ in each selector
$y_j$, and total degree at most $D$ in all ten amplitudes together.
Repeated univariate interpolation therefore leaves a nonzero
amplitude polynomial at some point of $T\times Y^{10}$.

The set $\Sigma_{10,D}$ detects every nonzero amplitude polynomial
of total degree at most $D$, over any characteristic.  To see this,
use the Newton basis
\[
 N_{\mathbf b}(\mathbf z)
  =\prod_{j=1}^{10}\prod_{h=0}^{b_j-1}(z_j-\theta_h),
 \qquad \sum_jb_j\le D.
\]
Its leading monomials are exactly those of total degree at most
$D$.  Evaluation at the node indexed by $\mathbf c$ vanishes
unless $c_j\ge b_j$ for every $j$.  Ordered by total degree,
the evaluation matrix is triangular with nonzero diagonal.
It is invertible, proving the claimed detection and (I.23).
Given the field nodes, computing the Lagrange polynomials and
enumerating the displayed family uses polynomially many field
operations in $n,D$, since its ten-pair parameter count is fixed.
Over a finite field, one may use an explicitly represented
extension large enough to supply the stated nodes.  Over
$\mathbb Q$ the nodes can be chosen as integers of polynomial
magnitude, giving rational coordinates of polynomial bit length.
\hfill$\square$

\paragraph{What this closes, and what it leaves open.}
This supplies a concrete positive answer for the width-three,
one-hidden-order case of I-G02, including arbitrarily many
rank-two cuts.  The rank-cut expression in I-19 equals
$f\prod_j\Delta_j$; the extra univariate factors can be absorbed
into the width-three layers, so the same theorem also hits those
particular coupled expressions with their increased degree bound.
The construction uses rank preservation directly and does not
derive a BIWA or a Hoza cocharacter decoder.
The general polynomial-in-$w$ gate I-G01 remains open in this
report.  In particular, I-08 requires widths growing with the
matrix dimension; this fixed-width result does not provide that
polynomial Noether-normalization conclusion.
